% FILE: artifacts/_common/style/type-roles.tex
% ==================================================================================================
% THE TYPE ROLES of the house style: the names that a document gives to its point sizes, and the rules
% that bind every type scale built on them. The VALUES are not here. Each document sets one size per
% role in its own type scale: the CCN 2026 poster in its tex/macros.tex, at true poster scale, and
% its supplement by re-declaring the figure roles at A4 scale. The math text helpers at the end of
% this file set words inside formulas; \brlab is pinned to the \notesize role.
%
% THE ROLES:
%   TEXT    \titlesize, \subtitlesize      poster title and its continuation
%           \headingsize                   block headings, authors, display formulae
%           \cardtitlesize                 card titles, heading kickers
%           \bodysize                      prose that wraps                    (PROSE leading)
%           \bodycompact                   body size, tight                    (COMPACT leading)
%           \notesize                      citations, metadata, tags, notes    -- TEXT FLOOR
%   FIGURE  \figcapsize                    figure captions, shared axis labels, colour keys
%           \figannotsize                  the figure's OBJECTS (vectors, patterns, matrices)
%           \figticksize                   subordinate marks: ticks, axes      -- FIGURE FLOOR
%   ICONS   \iconsm, \iconmd, \iconlg, \iconxl, \iconxxl -- FontAwesome glyphs in a disc
%
% TWO READABILITY FLOORS, read at different distances. Nothing may be set below its own, and each
% type scale states the two values at its own physical scale.
%   TEXT   -- anything read as language: prose, headings, captions, tags, notes, references, card
%             content, INCLUDING text drawn inside a TikZ card.
%   FIGURE -- annotations inside a diagram: axis ticks, symbol labels, marks.
% ICON GLYPHS are not text: not read as language, exempt from both floors, and sized to the disc that
% holds them, never to the type scale. Named so they cannot drift either.
%
% INSIDE A FIGURE, TWO LEVELS AND ONLY TWO:
%   \figannotsize -- the OBJECTS the figure is about: vectors, patterns, matrices.
%   \figticksize  -- everything SUBORDINATE to them: axis names, coefficients, ticks.
% A 30:24 ratio is a step; 36:24 reads as an accident rather than a hierarchy. \figcapsize is for
% WORDS about a figure -- its caption, a shared axis label, a colour key -- never for a symbol inside
% one. It shares \notesize's point size: a caption inside a figure and a note outside one are the same
% kind of small text, and one value keeps them from drifting apart.
%
% FONT SIZE IS ABSOLUTE INSIDE FIGURES, deliberately. A figure scales its COORDINATES (`scale=1.4')
% and never its fonts: no `transform shape', no \scalebox around a picture. Two consequences:
%   - Reusing a figure at another physical scale means RE-DECLARING these sizes (see supplement.tex).
%   - A label's CLEARANCE does not scale with it. Enlarging a picture moves a label further out in
%     coordinates while its glyph stays put, so it buys no room: raising a size forces a re-check of
%     every anchor in the figure, and collisions are fixed by MOVING THE ANCHOR, never by dropping
%     below the floor.
%
% THREE LEADINGS, NOT ONE:
%   PROSE   1.40x -- running text that wraps for several lines. Aerated, for reading at distance.
%   COMPACT 1.15x -- headings, captions, labels, notes: text that wraps to two or three lines and
%                    must stay a visual unit. A 1.4 leading pulls such a block apart.
%   FLAT    1.00x -- single-line nodes and icon glyphs, where leading is inert.
% When a gap is wanted that none of the three gives, ASK FOR THE GAP; do not invent a fourth ratio.
%
% Every step that can carry MATH needs a matching \DeclareMathSizes line in the same type scale, or
% its math silently drifts out of the system.

% ===== Math text helpers ==========================================================================
% In-math text label rendered in the math (serif) text font, so words match the symbols.
\newcommand{\mlab}[1]{\text{{\mathtextfont #1}}}
% An \underbracket's READING, pinned to \notesize INDEPENDENTLY of the formula's size: the gloss is not
% part of the statement, it is the words that READ the statement, and it must stay quiet. \mlab would
% follow the surrounding math size and grow with the formula. Use \brlab under any brace.
% NOTE: THE GLOSS, NOT THE ALGEBRA, SETS AN \underbracket's WIDTH -- it is the wider of the two.
\newcommand{\brlab}[1]{\text{{\mathtextfont\notesize #1}}}
% Inline math in a TikZ node that must read at prose size. A NAME ONLY: it must not transform anything.
% \DeclareMathSizes already reaches inside a TikZ node, so a \scalebox here would enlarge a second time.
\newcommand{\tmath}[1]{\ensuremath{#1}}
